% \begin{abstract}
% % Financial multi-agent systems built on large language models improve themselves exclusively through verbal reinforcement, reflecting on past trades in natural language, so risk-adjusted outcomes never update model parameters. When parameters \emph{are} updated, the target is directional accuracy, which induces reward hacking: hit-rate rises while cumulative return falls. We introduce \textsc{FinOPD}, which closes this loop through on-policy self-distillation conditioned on hindsight risk-adjusted PnL, Shapley-weighted credit assignment across specialist agents, and a non-parametric belief store. At the decision layer, Regime-Adaptive Signal Weighting and Edge-Gated Abstention restrict trading to market states where the system holds a measurable edge. Under strict post-cutoff evaluation on Dow-30 with an identical-protocol baseline suite, \textsc{FinOPD} attains SR\,=\,1.28 (vs.\ 0.79 for the strongest baseline) and reduces maximum drawdown 15--50$\times$.

% Financial multi-agent systems built on large language models increasingly decompose trading decisions into specialized analyst, risk, and portfolio-management roles, yet their self-improvement loops remain largely verbal: realized outcomes are summarized into later prompts without updating the policy that produced them. Direct parameter updates are also delicate, because optimizing directional accuracy can reward the wrong behavior: a model may predict more signs correctly while reducing cumulative return and worsening drawdown. We introduce \textsc{FinOPD}, a profit-anchored self-evolution framework that closes this loop with on-policy self-distillation conditioned on hindsight risk-adjusted PnL, Shapley-weighted credit assignment across specialist agents, and a non-parametric belief store for recurring market geometries. At the decision layer, Regime-Adaptive Signal Weighting and Edge-Gated Abstention restrict exposure to regimes where the system has a measurable edge. On a Dow-30 showcase set evaluated from January 2025 to May 2026 under a strict post-cutoff protocol, against an identical-protocol suite of trained time-series models and multi-trade agent baselines, \textsc{FinOPD} achieves the strongest portfolio-level risk-adjusted performance, reaching 60.9\% cumulative return, 1.41 Sharpe, and 3.30 Calmar, and remains consistently competitive across temporal and regime-conditioned analyses.

% \end{abstract}
% % 中文翻译：基于大语言模型的金融多智能体系统正日益把交易决策拆分给专门的分析、风控和组合管理角色，但它们的自我改进循环仍主要停留在语言层面：已实现结果被总结到后续提示词中，却不会更新产生这些决策的策略。直接更新参数也很敏感，因为优化方向准确率可能奖励错误行为：模型可能预测对更多涨跌方向，却降低累计收益并恶化回撤。我们提出 \textsc{FinOPD}，一个收益锚定的自进化框架，通过以事后风险调整 PnL 为条件的在线策略自蒸馏、跨专家智能体的 Shapley 加权信用分配，以及面向重复市场几何形态的非参数 belief 存储来闭合这一循环。在决策层，Regime-Adaptive Signal Weighting 与 Edge-Gated Abstention 将风险暴露限制在系统具有可测边际优势的 regime 中。在 2025 年 1 月至 2026 年 5 月从 Dow-30 中选取的 showcase 股票集合上的严格 post-cutoff 评估中，面对由训练好的时序模型和多笔交易智能体基线组成的同协议套件，\textsc{FinOPD} 取得最强的组合级风险调整表现，累计收益 60.9\%、Sharpe 1.41、Calmar 3.30，并在时间窗口和 regime 条件化分析中保持稳定竞争力。

\begin{abstract}
Financial multi-agent systems assemble rich evidence, yet their policies rarely learn directly from realized outcomes. Portfolio feedback is delayed, shared across agents, and must not leak hindsight into deployment. We present \textsc{FinOPD}, a temporally staged framework for learning from outcome-matured, current-policy trajectories. Before optimization, an LLM research process builds and freezes 157 executable alpha factors. During On-Policy Distillation (OPD), a point-in-time student generates trajectories; only after each horizon closes does a separate frozen teacher receive an estimated hindsight summary. Agent-weighted token JSD transfers supervision, while a token mismatch cap and reference-policy KL guard stabilize updates. Time-safe memory retains mature experience, and regime-aware selective execution maps evidence to cost-aware exposure. In a retrospective point-in-time backtest on four U.S. large-cap equities from January 2025 to May 2026, \textsc{FinOPD} ranks first among 13 methods by macro-averaged per-asset return, Sharpe, and Calmar (60.9\%, 1.41, and 3.30). We release synchronized-PnL evaluation and dependence-aware bootstrap code; these descriptive results do not constitute a live or contamination-free test.
% 中文翻译：金融多智能体系统能够汇集丰富证据，但其策略很少直接从已经实现的结果中学习。组合反馈具有延迟性，由多个智能体共同承担，而且部署过程中不得泄露事后信息。我们提出 FinOPD，一个从结果已经成熟的当前策略轨迹中学习的时间分阶段框架。优化开始前，LLM 研究流程构建并冻结 150+ 个可执行 Alpha 因子。On-Policy Distillation（OPD）期间，只使用时点信息的学生生成轨迹；只有在每个结果窗口闭合后，独立冻结教师才接收估计的事后摘要。按智能体加权的 token JSD 迁移监督信号，token 失配上限和参考策略 KL 约束共同稳定更新。时间安全记忆保留已经成熟的经验，状态感知的选择性执行则将证据映射为考虑成本的风险暴露。在 2025 年 1 月至 2026 年 5 月的四只美国大盘股上，采用相同成本、滑点和执行延迟时，FinOPD 在 13 个方法中取得最佳聚合收益、Sharpe 和 Calmar，累计收益达到 60.9\%，Sharpe 为 1.41，Calmar 为 3.30。
\end{abstract}
